\section{Simulation Verification}\label{sec:simulation}
A seven-degree-of-freedom nominal positive-definite joint-space model is used to evaluate numerical feasibility. The step size is $T_s=10^{-3}\,\mathrm{s}$ and the duration is $10\,\mathrm{s}$. A fifth-order polynomial reference is generated over $8\,\mathrm{s}$. The delay is $h=1\,\mathrm{s}$, with $R_h(\theta)=\sin^4(\pi(\theta-h)/h)$ on the second half-period. The gains are $K_p=\operatorname{diag}(25,25,25,22,22,20,20)$ and $K_d=\operatorname{diag}(10,10,10,9,9,8,8)$, with $K_0=[-K_p\ -K_d]$. Disturbances are $d_i=0.15\sin(0.7t+0.3i)+0.05\cos(1.3t+0.2i)$. Observer parameters are $T_o=T_c=1\,\mathrm{s}$, $\eta=0.3$, and $\sigma=0.4$.

\subsection{Results and Discussion}
The simulations show bounded base reaction motion, end-effector motion, joint tracking, and control torques under disturbance. The Gramian condition number is $5.796\times10^5$; consequently PDF and PD have nearly identical performance for this 14-state system. Without disturbance, terminal errors are about $10^{-9}$ rad. With disturbance, PD/PDF terminal error is approximately $1.4\times10^{-4}$ rad, whereas PDF+PTDO gives $3.327\times10^{-4}$ rad and a peak torque of $1920.6$ N m, indicating observer overcompensation under the selected discrete-time parameters. The result motivates larger $T_c$, parameter tuning, and saturation limits.

\begin{figure*}[t]
\centering
\subfloat[Free-floating trajectory\label{fig:free_floating_trajectory}]{\includegraphics[width=0.48\textwidth]{figures/free_floating_trajectory.pdf}}
\hfill
\subfloat[Base position\label{fig:base_position}]{\includegraphics[width=0.48\textwidth]{figures/base_position.pdf}}\\[1ex]
\subfloat[End-effector position\label{fig:end_effector_position}]{\includegraphics[width=0.48\textwidth]{figures/end_effector_position.pdf}}
\hfill
\subfloat[Joint posture\label{fig:joint_posture}]{\includegraphics[width=0.48\textwidth]{figures/joint_posture.pdf}}
\caption{Free-floating base, end-effector, and joint trajectories. Each panel is exported independently by MATLAB and assembled using LaTeX.}
\label{fig:motion}
\end{figure*}

\begin{figure*}[t]
\centering
\subfloat[Joint position tracking\label{fig:joint_position}]{%
\includegraphics[width=0.48\textwidth]{figures/joint_position_tracking.pdf}}
\hfill
\subfloat[Joint velocity tracking\label{fig:joint_velocity}]{%
\includegraphics[width=0.48\textwidth]{figures/joint_velocity_tracking.pdf}}\\[1ex]
\subfloat[Position and velocity tracking errors\label{fig:tracking_errors}]{%
\includegraphics[width=0.48\textwidth]{figures/tracking_errors.pdf}}
\hfill
\subfloat[Control torque\label{fig:control_torque}]{%
\includegraphics[width=0.48\textwidth]{figures/control_torque.pdf}}
\caption{Joint tracking results: desired versus actual joint position and velocity, tracking errors, and control torque.}
\label{fig:joint_summary}
\end{figure*}

\begin{table}[t]\centering\caption{Simulation metrics}\label{tab:sim_metrics}\begin{tabular}{cccc}\toprule Controller and condition & $\|e(T)\|$ [rad] & $\mathrm{RMS}(e)$ [rad] & $\max|\tau_i|$ [N m]\\\midrule PD, no disturbance & $1.333\times10^{-9}$ & $5.039\times10^{-10}$ & 258.9\\ PDF, no disturbance & $1.333\times10^{-9}$ & $5.039\times10^{-10}$ & 258.9\\ PD, disturbance & $1.373\times10^{-4}$ & $5.188\times10^{-5}$ & 258.9\\ PDF, disturbance & $1.373\times10^{-4}$ & $5.188\times10^{-5}$ & 258.9\\ PDF+PTDO, disturbance & $3.327\times10^{-4}$ & $1.258\times10^{-4}$ & 1920.6\\\bottomrule\end{tabular}\end{table}

\begin{figure*}[t]
\centering
\subfloat[True and estimated disturbance\label{fig:observer_estimate}]{\includegraphics[width=0.48\textwidth]{figures/observer_estimates.pdf}}
\hfill
\subfloat[Observer error\label{fig:observer_error}]{\includegraphics[width=0.48\textwidth]{figures/observer_error.pdf}}
\caption{Disturbance-observer response.}
\label{fig:observer}
\end{figure*}